\documentclass[preprint,12pt,authoryear]{elsarticle}
\usepackage[utf8]{inputenc}
\usepackage{amsmath,amssymb}
\usepackage{booktabs}
\usepackage{graphicx}
\usepackage{setspace}
\usepackage[hidelinks]{hyperref}
\usepackage{enumitem}
\usepackage{lineno}

\journal{Preventive Medicine}

\begin{document}

\begin{frontmatter}

\title{Waiting Time as a Behavioral Barrier to Vaccine Uptake: The Role of Institutional Trust and Medical Vulnerability}

\author[1]{Yichao Jin}
\author[1]{Dohyeong Kim\corref{cor1}}
\address[1]{School of Economic, Political and Policy Sciences, The University of Texas at Dallas, Richardson, TX, USA}
\cortext[cor1]{Corresponding author. Email: dohyeong.kim@utdallas.edu}

\begin{abstract}
\noindent \textbf{Objective:} To quantify waiting time as a behavioral barrier to vaccination and evaluate how institutional trust and medical vulnerability modify the perceived burden of delay in preventive service delivery.\\
\textbf{Methods:} A discrete choice experiment (N=1,027) in Wuhan, China, assessed preferences for vaccine efficacy, safety, and wait times. Mixed Logit and Latent Class models estimated the marginal willingness-to-accept (MWTA) for delay and the perceived delay burden (the subjective-to-objective wait time ratio).\\
\textbf{Results:} Waiting time was a primary determinant of uptake. On average, respondents required 47 RMB per month to offset delay disutility. Institutional trust significantly moderated preferences; low-trust individuals exhibited 2.2 times greater delay sensitivity than high-trust peers. For medically vulnerable groups, the perceived delay burden reached 1.5, indicating that administrative friction disproportionately affects those at highest biological risk.\\
\textbf{Conclusions:} Service delays impose a significant behavioral cost that may undermine vaccination programs. Public health strategies should prioritize operational transparency in wait-time communication and implement streamlined delivery for vulnerable populations. Enhancing institutional trust is an important intervention to reduce the perceived burden of delay and ensure rapid, equitable access to preventive care.
\end{abstract}

\begin{keyword}
Vaccine uptake \sep Waiting time \sep Institutional trust \sep Public health policy \sep Behavioral barriers \sep Health equity
\end{keyword}

\vspace{0.5cm}
\noindent \textbf{Word Count:} Abstract: 153 words | Main Text: $\sim$2,600 words | References: 30 | Tables/Figures: 4

\end{frontmatter}

\doublespacing
%\linenumbers

\section{Introduction}
Vaccination remains a cornerstone of infectious disease prevention and a primary objective of global health policy. However, the success of immunization programs depends not only on final coverage levels but also on the critical dimension of uptake timing \citep{Chapman2005, Chapman2010}. In the context of infectious disease outbreaks, every day of delay represents a window of biological vulnerability that can amplify transmission risk and strain public health resources \citep{Castioni2024, Shankar2024}. While logistical obstacles such as supply chain management are well-documented, the behavioral barriers imposed by service delays and waiting times remain less understood in implementation research.

Current health promotion strategies often prioritize overcoming vaccine hesitancy through educational campaigns focused on clinical efficacy and safety profiles \citep{MacDonald2015, Anderson2023, Betsch2018}. While essential, these approaches may overlook the administrative friction embedded within delivery systems. From a behavioral perspective, waiting for a vaccine is not a neutral logistical event; it is an intertemporal trade-off where individuals must weigh immediate opportunity costs against the future benefit of acquired immunity \citep{Frederick2002, Attema2025}. If the perceived burden of waiting is high, even highly efficacious vaccines may face suboptimal uptake during critical periods, leading to preventable morbidity.

The impact of service delays may not be uniform across the population. Evidence suggests that institutional trust—defined here as confidence in health authorities and government systems—may function as a critical moderator of health-seeking behavior. Furthermore, individuals with underlying medical conditions or higher biological risk may perceive the costs of delay more acutely. Understanding how institutional trust and medical vulnerability shape responses to administrative friction is essential for designing equitable vaccination programs that do not inadvertently penalize the most at-risk segments of the population.

This study utilizes a large-scale discrete choice experiment (N=1,027) conducted in Wuhan, China, to quantify waiting time as a behavioral barrier to vaccination. We aim to determine how this barrier competes with clinical attributes such as efficacy and safety and to evaluate the roles of institutional trust and medical risk in shaping these trade-offs. By reframing service delivery speed as a behavioral determinant of health, we provide an evidence-based framework to help public health officials reduce administrative friction, enhance communication transparency, and improve the overall efficiency of preventive care systems.

\section{Methods}
\subsection{Study Population and Policy Context}
We conducted a cross-sectional study in Wuhan, China, a metropolitan area that provides a unique context for investigating health-related time preferences due to the high historical and social salience of infectious disease management. A total of 1,027 adults (18 years and older) participated in the study. The sampling frame utilized a multi-stage stratified approach, ensuring representativeness across diverse socioeconomic backgrounds, age groups, and district-level population densities.

The survey was implemented through a secure digital platform in [AUTHOR CHECK: add exact survey dates, e.g., October 2024]. This period was characterized by a transition in public health policy where the focus shifted from emergency mass vaccination to long-term sustainable prevention. Participants provided informed consent, and the study design was approved by the institutional review board [AUTHOR CHECK: add IRB institution and approval number]. The survey instrument was pre-tested through pilot interviews to ensure that the descriptions of vaccine attributes were easily understood by individuals without medical training.

\subsection{Experimental Design: Quantifying Implementation Barriers}
The core of our methodology is a Discrete Choice Experiment (DCE), which is widely used in health services research to elicit latent preferences by forcing participants to make trade-offs between competing service attributes. The DCE format is particularly effective for identifying the "relative importance" of non-clinical factors like waiting time.

Each participant was presented with 12 choice tasks. In each task, they were asked to choose between three hypothetical vaccination programs and an opt-out (no vaccination) option. The attributes and levels were carefully selected to reflect the primary levers available to public health administrators (see Table \ref{tab:attributes}):
\begin{enumerate}[label=(\arabic*)]
    \item \textbf{Vaccine Efficacy (50\%, 70\%, 90\%)}: The primary clinical outcome of the intervention.
    \item \textbf{Side Effects (Mild, Moderate, Severe)}: The perceived safety cost associated with uptake.
    \item \textbf{Cash Incentives (0, 200, 500 RMB)}: A common policy tool used to offset the costs of health-seeking behavior \citep{Bechtel2023}.
    \item \textbf{Waiting Time (Immediate, 1 month, 3 months, 6 months)}: The primary administrative barrier, representing the time from the decision to vaccinate until the protection is active.
\end{enumerate}

\begin{table}[ht]
\centering
\caption{Attributes and Levels in the Discrete Choice Experiment}
\label{tab:attributes}
\begin{tabular}{ll}
\toprule
Attribute & Levels \\
\midrule
Vaccine Efficacy & 50\%, 70\%, 90\% \\
Side Effects & Mild, Moderate, Severe \\
Cash Incentives & 0, 200, 500 RMB \\
Waiting Time & Immediate, 1, 3, 6 months \\
\bottomrule
\end{tabular}
\end{table}

We employed a D-efficient experimental design (using Ngene software) to maximize the statistical power of the choice data. The conceptual relationship between these attributes and the decision to vaccinate is illustrated in Figure \ref{fig:framework_supp} (see Supplementary Material).

\subsection{Econometric Specification and Implementation Metrics}
We analyzed the choice data using a Random Utility Framework \citep{mcfadden2001, Train2009, Hole2007}. Specifically, we utilized Mixed Logit (MXL) models to identify both the average population preference and the degree of taste heterogeneity. To further explore how these preferences vary across distinct population segments, we employed Latent Class Analysis (LCA), which identifies unobserved subgroups (classes) with similar choice patterns related to socioeconomic characteristics and health beliefs.

To translate these findings into actionable evidence for preventive medicine, we estimated two primary implementation metrics:
\begin{enumerate}
    \item \textbf{Marginal Willingness-to-Accept (MWTA) for Delay}: This measures the monetary value (in RMB) required to compensate a participant for one additional month of waiting. It serves as a direct proxy for the behavioral ``friction'' or disutility imposed by the delivery system.
    \item \textbf{Perceived Delay Burden (PDB)}: We estimated the ratio of the behavioral cost of delay to its objective chronological duration. A ratio significantly greater than 1.0 indicates that participants perceive a disproportional burden, where the wait is felt more acutely than its actual length.
\end{enumerate}
All models were estimated using maximum simulated likelihood with cluster-robust standard errors and 1,000 Halton draws to ensure robust results.

\section{Results}
\subsection{Waiting Time as a Dominant Barrier to Prevention}
The Mixed Logit results confirm that waiting time is a primary determinant of vaccination uptake ($p < 0.01$; see Table \ref{tab:results_main} for full model estimates). Across the population, the disutility associated with a 3-month wait was found to be statistically equivalent to a significant (approx. 20\%) reduction in vaccine efficacy. This finding suggests important considerations for health promotion: it suggests that the public health benefits of high-efficacy vaccines may be reduced if the delivery system cannot provide ``fast'' protection.

The average respondent required 47 RMB per month of delay to maintain the same probability of uptake. This suggests that for every month of administrative delay, the ``perceived value'' of the vaccination program diminishes, making individuals more likely to opt out or postpone uptake, thereby increasing the risk of community-level transmission.

\begin{table}[ht]
\centering
\caption{Mixed Logit Model Estimates: Behavioral Barriers to Vaccination}
\label{tab:results_main}
\begin{tabular}{lrr}
\toprule
Attribute & Coefficient (S.E.) & $z$-stat \\
\midrule
Vaccine Efficacy (per \%) & 0.042 (0.003)*** & 14.00 \\
Side Effects (Severe vs. Mild) & -1.850 (0.120)*** & -15.42 \\
Cash Incentives (per 100 RMB) & 0.150 (0.015)*** & 10.00 \\
Waiting Time (per month) & -0.075 (0.008)*** & -9.38 \\
Institutional Trust (Interaction) & 0.055 (0.012)*** & 4.58 \\
\midrule
Log-likelihood & -8,450.2 & \\
McFadden $R^2$ & 0.285 & \\
Observations (N) & 1,027 & \\
\bottomrule
\multicolumn{3}{l}{*** $p < 0.01$, ** $p < 0.05$, * $p < 0.1$}
\end{tabular}
\end{table}

\subsection{The Moderating Role of Institutional Trust}
One of the most significant findings for public health policymaking is the interaction between institutional trust and time preferences. We partitioned the sample based on self-reported trust in health authorities and government institutions. 

As shown in Figure \ref{fig:trust_nexus}, trust acts as a critical ``delivery lubricant.'' Participants with high institutional trust were significantly more patient and more willing to tolerate logistical delays (MWTA = 28.7 RMB). In contrast, for individuals with low trust, the behavioral cost of waiting nearly doubled (MWTA = 62.3 RMB).

\begin{figure}[ht]
\centering
\includegraphics[width=0.8\textwidth]{fig2.png}
\caption{The trust-time nexus: Marginal Willingness-to-Accept (MWTA) for delay by institutional trust level.}
\label{fig:trust_nexus}
\end{figure}

This trust gap indicates that in the absence of institutional confidence, the ``shadow of the future'' shortens, and the perceived risk of waiting becomes a dominant barrier to participation in preventive programs.

\subsection{Subgroup Analysis: A Social Determinants Lens}
To further align our findings with the social determinants of health (SDoH) perspective, we explored how choice behavior varies across age and income strata using Latent Class Analysis. Our results revealed significant heterogeneity in how different population segments prioritize vaccine attributes. Elderly respondents (age $>$ 60) exhibited the highest sensitivity to side effects but were more willing to tolerate waiting times if efficacy was high. Conversely, low-income participants showed a significantly higher Marginal Willingness-to-Accept (MWTA) for cash incentives, confirming that financial barriers remain a primary obstacle to equitable vaccine coverage.

Crucially, the interaction between income and waiting time revealed a ``double burden'' for the socioeconomically disadvantaged. Individuals in the lowest income quartile reported a significantly higher Perceived Delay Burden (PDB = 1.50). This finding supports the necessity of localized, high-speed delivery models that minimize the ``time cost'' of prevention for low-income communities.

\subsection{Wald Tests and Structural Robustness}
We conducted a series of Wald tests to verify the statistical significance of the differences in PDB across subgroups. The null hypothesis that the PDB for the medically vulnerable group (1.50) is equal to the non-vulnerable group (1.18) was rejected at the 1\% level ($\chi^2 = 14.32, p < 0.001$). Similarly, the difference in MWTA for delay between high-trust and low-trust respondents (33.6 RMB difference) was found to be highly significant ($\chi^2 = 9.87, p < 0.01$). These results provide robust statistical support for the claim that institutional trust and medical risk fundamentally alter the psychological ``price'' of waiting for preventive services.

\section{Discussion}
Our findings demonstrate that waiting time functions as a significant behavioral barrier to vaccination uptake, comparable in perceived disutility to substantial reductions in clinical efficacy. This suggests that for many individuals, the time-to-protection is as critical a determinant of health-seeking behavior as the vaccine's safety or performance profile. By quantifying these intertemporal trade-offs, this study suggests the potential benefit of incorporating administrative speed as a relevant component of vaccination program design.

A key contribution of this research is the identification of the moderating roles of institutional trust and medical vulnerability. The observation that low-trust individuals exhibit significantly higher sensitivity to service delays suggests that institutional confidence serves as a stabilizing factor during delivery uncertainty. In low-trust environments, administrative friction may exacerbate suspicion and disengagement, potentially creating a negative feedback loop that depresses uptake. Furthermore, the heightened delay burden perceived by medically vulnerable and socioeconomically disadvantaged groups indicates that administrative delays may act as a regressive barrier, disproportionately penalizing those at the highest biological and financial risk.

These results have several practical implications for public health implementation. First, increasing operational transparency regarding vaccine availability and wait times may reduce the uncertainty that often inflates the perceived burden of delay. Second, our findings support the development of tailored delivery models, such as ``fast-track'' protocols for high-risk populations. By prioritizing access for those who perceive the costs of waiting most acutely, health authorities can improve both the efficiency and the equity of preventive programs. Finally, rebuilding institutional trust through consistent and transparent communication is likely important for the success of long-term immunization strategies.

This study has certain limitations. First, as a discrete choice experiment, the results are based on stated preferences in a hypothetical scenario rather than revealed behavior from actual vaccination programs. While such methods are validated for eliciting latent trade-offs, future research using real-world registry data or randomized trials is needed to confirm the impact of wait-time interventions on actual coverage rates. Second, the study was conducted in a specific metropolitan context (Wuhan), and the generalizability of these ``time-trust'' dynamics to other health systems with different institutional structures remains to be explored.

In conclusion, service delays represent a meaningful behavioral barrier that can undermine the effectiveness of disease prevention efforts. By addressing administrative friction and institutional trust as modifiable determinants of uptake, public health officials can design more responsive vaccination delivery systems. Prioritizing timely and equitable access for vulnerable populations is essential to maximize the public health impact of both existing and future immunization programs.

\section*{Conflict of Interest Statement}
The authors declare that they have no known competing financial interests or personal relationships that could have appeared to influence the work reported in this paper.

\section*{Acknowledgments}
The authors would like to thank the study participants in Wuhan for their time and contribution to this research.

\section*{References}
\bibliographystyle{elsarticle-harv}
\bibliography{refs}

\newpage
\appendix
\section{Supplementary Material}
\renewcommand{\thefigure}{S\arabic{figure}}
\setcounter{figure}{0}

\begin{figure}[ht]
\centering
\includegraphics[width=0.8\textwidth]{fig1.png}
\caption{Conceptual framework of vaccination timing as a behavioral barrier.}
\label{fig:framework_supp}
\end{figure}

\end{document}
